\subsection{Fundamentals of travel-time tomography}

Travel-time tomography (TTT) \cite{bormann_seismic_2012,stefanov_travel_2019} is one of the simplest methods for SSF estimation. It uses the difference in arrival time of the acoustics signals at the sensors, requiring a spatial distribution of sources and receiver(s) for sufficient spatial resolution.

In general, the SSF is discretized in layers or cells, where each behaves as having a constant sound speed. Then, the wavepath's length in each layer is calculated, resulting in a system of equations that is used to estimate the velocity in each layer. It can also work with assumed functions for the SSF, ditching the layering assumption, but increasing the problem's complexity.

Its simpleness, robustness and computational efficiency are its main advantages, the later allowing real-time implementation. However, it depends on a precise positioning and syncing of the devices, and its discretization assumption results in limited resolution.

\begin{figure}[H]
	\centering
	\begin{tikzpicture}[scale=1.5]
		\coordinate (tl) at (-3, 2);
		\coordinate (tr) at (3, 2);
		\coordinate (bl) at (-3, -2);
		\coordinate (br) at (3, -2);
		\coordinate (cc) at (0,0);
		\coordinate (tc) at (0,2);
		\coordinate (bc) at (0,-2);
		\coordinate (cl) at (-3, 0);
		\coordinate (cr) at (3,0);
		\coordinate (source) at ($(bl)!0.95!(br)$);
		\coordinate (receiver) at ($(bl)!0.9!(tl)$);
		\draw[gray,fill=none] (tl) rectangle (br);
		\draw[->-,thick] (source) to[out=120,in=-18] (receiver) node[pos=0.5,anchor=center,black] {$\bvr[v]$};
		%		\filldraw[blue!50!black] (receiver) rectangle (2,2);
		%		\filldraw[green!50!black] (source) circle (3pt);
		\node[square,fill=blue!70!black] () at (receiver) {};
		\node[triangle,fill=green!70!black] () at (source) {};
		
	\end{tikzpicture}
	\caption{Wave path on heterogeneous and continuously variable environment.}
	\label{fig:sec2:wavepath_heterogeneous-variable_environment}
\end{figure}

Consider a wave that follows the path from \cref{fig:sec2:wavepath_heterogeneous-variable_environment}, this being represented by the curve $\bvr[v]$. Assuming that the speed of sound $c(\bvp)$ is a function of space ($\bvp = [x,y,z]$), then the travel time along the curve $\bvr[v]$ is
\begin{equation}
	t_{\bvr[v]} = \int_{\bvr[v]} \frac{1}{c(\bvp)} \dd\bvp.
\end{equation}

The objective is to, given a set of $N$ measures $t_{\bvr[v]{i}}$, these representing different wave paths, estimate the sound speed field $c(\bvp)$.

Different approaches can be taken to solve this problem, including to discretize space, approximate the SSF as polynomial through Spline or least-squares fitting, among others. This technique can also be employed in 1D, 2D or 3D situations, and acquiring spatial diversity through either (or both) multi-source or multi-receiver means.

Most models estimate the slowness (inverse of speed) field. Note that SSF can both mean \quote{sound speed field} and \quote{sound slowness field}. Since both quantities are the reciprocal of one another, the SSF acronym will be indiscreetly used for both.

\subsection{Standard discrete travel-time tomography approach}
\label{subsec:sec2:standard_ttt_approach}

All of the cited papers took an approach similar to the one established below. Differences may arise from ignoring ray bending, type of regularization in the minimization process, amount and positioning of sources and receivers, final purpose for the TTT, or some other particularization or assumption.

In these regards, the most complete papers are the ones by Ali \cite{ali_opensource_2019}, Phillips \cite{phillips_traveltime_1991}, and Hormati \cite{hormati_robust_2010}, where they consider regularization, ray bending, and a large amount of sources and receivers. We now will lay the common ground among all the presented papers, the median modeling considerations and formulations, as well as some important distinctions and particularities that are shared between some of them. Note that the framework presented here isn't identical to the ones in the cited papers \cite{aki_determination_1977,aki_determination_1976,ali_opensource_2019,hormati_robust_2010,jeong_investigating_2024,hellmann_borehole_2023,phillips_traveltime_1991,zhang_nonlinear_1998,dines_computerized_1979}, and only highlights the similarities and mostly-shared ground between them.

Let's assume the environment to be modeled is divided into an $\sz{N_1}{N_2}$ grid of cells (horizontal $\mtimes$ vertical), with a total of $N = N_1\cdot N_2$ cells. A total of $M$ rays are produced, between any number of sources and receivers. We denote $\bvs \in \re^{\sz{N}{1}}$ positive vector representing the slowness in each cell, and $\bvR{m}(\bvs) \in \re^{\sz{N}{1}}$ vector containing the length traversed in each cell, by the $m$-th ray. The length-vector $\bvR{m}(\bvs)$ is commonly obtained externally via a fast-marching algorithm or another solver for the Eikonal equation, considering Fermat's principle of least time, and Snell's law $\bvR{m}(\bvs)$. $\bvR{m}(\bvs)$ can also be obtained directly if one considers an SSF with small perturbations, where ray-bending effects can be ignored, and the path is assumed linear between source and receiver. With this, the travel-time for the $m$-th ray is $t_m$, being given by
\begin{equation}
	t_m(\bvs) = \tr{\bvr{m}}(\bvs) \bvs
\end{equation}

Now we let $\bvt(\bvs) \in \re^{\sz{M}{1}}$ vector of travel-times for the current model, for all $M$ considered rays; and $\bvR(\bvs) \in \re^{\sz{M}{N}}$ matrix, where the $m$-th row is $\tr{j_m}(\bvs)$. With this, our minimization problem becomes minimizing the error $E(\bvs)$, this being translated into
\begin{equation}
	\opt{\bvs} = \arg\min_{\bvs} E(\bvs) = \norm{\bvR(\bvs) \bvs - \bvt{\obs}}^2
	\label{eq:sec2:def_minimization_bvz}
\end{equation}
where $\bvt{\obs}$ is the observed travel-time for each ray.

Note that this problem is non-linear and recursive: the solution to the ray paths, given by $\bvR(\bvs)$, depends on the slowness field $\bvs$. Meanwhile, the field model depends on the achieved solution for the ray paths, $\bvR(\bvs)$. Given the non-linearity of the problem, this search is usually done iteratively, following
\begin{subequations}
	\begin{gather}
		\bvs{i+1} = \arg\min_{\bvs{i+1}} E(\bvs{i+1}) = \norm{\bvR(\bvs{i}) \bvs{i+1} - \bvt{\obs}}^2 
		\label{subeq:sec2:iterative_minimization:eq1}\\
		\nabla E(\bvs{i+1}) = 2\tr{\bvR}(\bvs{i}) \pts{\bvR(\bvs{i}) \bvs{i+1} - \bvt{\obs}}
		\label{subeq:sec2:iterative_minimization:eq2}
	\end{gather}
	\label{eqs:sec2:iterative_minimization}
\end{subequations}

Usually, this problem is ill-posed, and requires extra conditions for a unique solution. The most common is regularization, which leads to the problem
\begin{subequations}
	\begin{gather}
		\bvs{i+1} = \arg\min_{\bvs{i+1}} E(\bvs{i+1}) = \norm{\bvR{i} \bvs{i+1} - \bvt{\obs}}^2 + P_\alpha(\bvs{i+1})\\
		\nabla E(\bvs{i+1}) = 2\bvR[T]{i} \pts{\bvR{i} \bvs{i+1} - \bvt{\obs}} + \nabla P_{\alpha}(\bvs{i+1})
	\end{gather}
	\label{eqs:sec2:iterative_minimization_Ezi}
\end{subequations}
with $\alpha$ being a controlling parameter; and, assuming that the regularization function can be given by $P_{\alpha}(\bvs{i+1}) = \alpha^2\norm{\bvD\bvs{i+1}}^2$ where $\bvD$ is a square matrix, its solution is
\begin{equation}
	\bvG{i} = \inv*{\tr{\bvR{i}} \bvR{i} + \alpha^2 \tr{\bvD}\bvD} \bvR[T]{i}
	\label{eq:sec2:solution_regularization-problem_simple}
\end{equation}
for this inversion step to be possible, it is sufficient that $\bvD$ is a full-rank matrix\footnote{\label{footnote:positive-definite_inversion}By definition, $\bvR[T]{i} \bvR{i}$ is positive-semidefinite, $\tr{\bvD}\bvD$ is positive-definite if it is full-rank square matrix, and the sum of a positive-semidefinite and a positive-definite is always positive-definite, this being a sufficient condition for invertibility.}.

The regularization matrix $\bvD$ can be chosen to improve a range of metrics. Choosing $\bvD = \Id$ \cite{phillips_traveltime_1991} results in the minimization also considering the root mean-square (RMS) of the slowness vector; this is equivalent to considering the presence in white noise in the measurements, and trading model accurateness for noise rejection. Choosing $\bvD z \approx \nabla^2 s(x,y)$, (discretely) approximating the Laplacian of the (continuous) slowness field \cite{zhang_nonlinear_1998,ali_opensource_2019}, resulting in a minimization that penalizes the Laplacian of the achieved slowness field; that is, a solution that also considers the curvature/jumps in the SSF, searching for a smoother field.

Another interesting tool that is seen in the literature \cite{phillips_traveltime_1991,meyerholtz_convolutional_1989} is to employ a blur on the gradient. That is, we define that
\begin{equation}
	E(\bvs{i+1}) = \norm{\bvR[d]{i} \bvs[d]{i+1} - \bvt{\obs}}^2 + P_\alpha(\bvs[d]{i+1})
\end{equation}
where $\bvR[d]{i} = \bvR{i}\bvB{i}$ and $\bvs[d]{i+1} = \inv{\bvB{i}} \bvs{i+1}$, and $\bvB{i}$ is the $i$-th iteration's blurring matrix (this being based on a Gaussian blur, or any other averaging kernel, that decreases in radius at each iteration). In this scenario, the solution $\bvs{i+1}$ is
\begin{equation}
	\bvG{i} = \bvB{i} \inv*{\tr{\bvB{i}} \tr{\bvR{i}} \bvR{i} \bvB{i} + \alpha^2 \tr{\bvD}\bvD} \tr{\bvB{i}} \bvR[T]{i}
\end{equation}

To achieve this, the regularization function has to be adapted, such that $P_{\alpha}(\bvs{i+1}) = \alpha^2 \norm{\bvD \inv{\bvB{i}} \bvs{i+1}}^2$; this means the regularization is applied on the ``reverse blurred'' signal. As stated, the blurring matrix $\bvB{i}$ depends on the iterator $i$. Usually, the blur radius is decreased for increasing $i$, aiming at gradually finding a more refined and accurate (less blurry) solution.

This convolutional averaging (also called ray broadening \cite{meyerholtz_convolutional_1989,phillips_traveltime_1991}) can help find solutions by scattering the rays and expanding the paths taken to neighboring cells, thus dealing with the problem of some cells not being traversed by any rays.